\section*{Capítulo \ref{ch:b3:supramolecular} --- Química supramolecular}

\begin{solution}{exo:b3:supramolecular:1}
Íon--dipolo; empilhamento $\pi$; cátion--$\pi$; \omterm{def:b3:supramolecular:interactions}{ligação de halogênio} (de I para N); três ligações de hidrogênio.
\end{solution}

\begin{solution}{exo:b3:supramolecular:2}
$\Delta G^\circ = -RT\ln K = -8.314 \times 298 \times \ln(\num{1.0e4}) = \qty{-22.8}{kJ/mol}$.
\end{solution}

\begin{solution}{exo:b3:supramolecular:3}
$K_{\ce{K+}}/K_{\ce{Na+}} = 10^{1.7} = 50$; a diferença é $RT\ln 50 = \qty{9.7}{kJ/mol}$ a favor do potássio.
\end{solution}

\begin{solution}{exo:b3:supramolecular:4}
$n/(n + m) = 1/3$: um \omterm{def:b3:supramolecular:supramolecular}{hospedeiro} para dois \omterm{def:b3:supramolecular:supramolecular}{hóspedes}, $\mathrm{HG_2}$.
\end{solution}

\begin{solution}{exo:b3:supramolecular:5}
$H_0 + G_0 + 1/K = \qty{4.667e-3}{mol/L}$; $[\mathrm{HG}] = (4.667 - \sqrt{4.667^2 - 4 \times 2.0 \times 2.0})/2 \times 10^{-3} = \qty{1.13e-3}{mol/L}$: fração
ligada 0.566, $\delta = 3.560 + 0.160 \times 0.566 = \qty{3.651}{ppm}$.
\end{solution}

\begin{solution}{exo:b3:supramolecular:6}
$\log K = 18.3 - 8.6 = 9.7$, $K = \num{5e9}$; $\Delta G^\circ = -RT\ln K = \qty{-55}{kJ/mol}$. As seis ligações Ni--N são do
mesmo tipo dos dois lados; o número de moléculas livres passa de quatro para
sete: o ganho de entropia translacional impulsiona a reação.
\end{solution}

\begin{solution}{exo:b3:supramolecular:7}
$\Delta G^\circ = -RT\ln(\num{2.0e5}) = \qty{-30.2}{kJ/mol}$; $T\Delta S^\circ = \Delta H^\circ - \Delta G^\circ = -40 + 30.2 = \qty{-9.8}{kJ/mol}$;
$\Delta S^\circ = \qty{-33}{J.K^{-1}.mol^{-1}}$: ligação impulsionada pela entalpia, em parte paga em entropia.
\end{solution}

\begin{solution}{exo:b3:supramolecular:8}
O cobre(I) liga duas unidades bidentadas de fenantrolina e, sendo tetraédrico, as mantém
em ângulo reto, uma enfiada no espaço onde o anel da outra vai
se fechar. Fechar as duas cadeias em anéis as entrelaça. Um grande excesso de
cianeto, que liga o cobre(I) muito mais fortemente, remove o metal e deixa
o \omterm{def:b3:supramolecular:interlocked}{catenano} livre.
\end{solution}

\begin{solution}{exo:b3:supramolecular:9}
$k_c = \pi\Delta\nu/\sqrt2 = \pi \times 120/1.414 = \qty{267}{s^{-1}}$. Eyring:
$\Delta G^\ddagger = RT\ln(k_BT/hk_c) = 8.314 \times 300 \times \ln(\num{6.25e12}/267) = \qty{59.6}{kJ/mol}$.
\end{solution}

\begin{solution}{exo:b3:supramolecular:10}
Com excesso de fármaco sólido, o fármaco livre fica fixado em $S_0$. Então $[\mathrm{DC}] = KS_0[\mathrm C]$ e
$C_t = [\mathrm C] + [\mathrm{DC}] = [\mathrm C](1 + KS_0)$, logo $[\mathrm{DC}] = KS_0C_t/(1 + KS_0)$ e a solubilidade total é $S_0 + [\mathrm{DC}]$. Aqui,
$KS_0 = 0.050$: $S = 0.10 + 0.050 \times 10/1.050 = \qty{0.58}{mmol/L}$, quase seis vezes o valor intrínseco.
\end{solution}

\begin{solution}{exo:b3:supramolecular:11}
Quando $1/K \ll H_0, G_0$, $[\mathrm{HG}] \approx \bigl(H_0 + G_0 - |H_0 - G_0|\bigr)/2 = \min(H_0, G_0)$: a fração ligada é $G_0/H_0$ até um
equivalente, depois 1. A curva deixa de depender de $K$: qualquer valor grande o bastante
dá a mesma quebra nítida dentro do ruído, de modo que os dados mostram apenas que $K$ excede
um certo mínimo.
\end{solution}

\begin{solution}{exo:b3:supramolecular:12}
$k_{\mathrm{off}} = k_{\mathrm{on}}/K = \qty{1e9}{L.mol^{-1}.s^{-1}}/\qty{1500}{L.mol^{-1}} = \qty{6.7e5}{s^{-1}}$. A diferença de deslocamento
$\qty{0.160}{ppm}$ corresponde a \qty{64}{Hz} a \qty{400}{MHz}; a coalescência exigiria $k \approx \pi \times 64/\sqrt2 = \qty{140}{s^{-1}}$. A troca é
milhares de vezes mais rápida: um único sinal médio.
\end{solution}

\begin{solution}{pb:b3:supramolecular:1}
\textbf{1.} Um anel de seis unidades $\ce{-CH2CH2O-}$; no complexo, o $\ce{K+}$ fica no centro, ligado pelos
seis oxigênios voltados para dentro, com os grupos $\ce{CH2}$ do lado de fora.

\textbf{2.} $\ce{K+}$: $-RT\ln 10 \times 6.1 = \qty{-34.8}{kJ/mol}$; $\ce{Na+}$: \qty{-25.1}{kJ/mol}.

\textbf{3.} $10^{6.1 - 4.4} = 50$.

\textbf{4.} O $\ce{K+}$, maior (\qty{138}{pm}), alcança os seis oxigênios do anel ao mesmo tempo; o
$\ce{Na+}$, menor (\qty{102}{pm}), não consegue, a menos que o anel se dobre, o que custa tensão.

\textbf{5.} Seu exterior é uma superfície de hidrocarboneto e sua carga está enterrada; o $\ce{KCl}$
precisaria ter seus íons solvatados, o que o tolueno não consegue fazer.

\textbf{6.} A água solvata o $\ce{K+}$ muito melhor que o metanol, e o íon precisa perder essa
água para entrar no anel.

\textbf{7.} $\ce{K+}(\text{aq}) + \ce{MnO4-}(\text{aq}) + \mathrm L(\text{org}) \rightleftharpoons \ce{[KL]+MnO4-}(\text{org})$.

\textbf{8.} $[\ce{K+}]_{\mathrm{aq}}$ permanece \qty{0.10}{mol/L}: $y = \num{1.0e5} \times 0.10 \times (\num{1.0e-3} - y)^2 = \num{1.0e4}(\num{1.0e-3} - y)^2$.

\textbf{9.} Com $u = \num{1.0e-3} - y$: $\num{1.0e4}u^2 + u - \num{1.0e-3} = 0$, $u = \qty{2.70e-4}{mol/L}$, $y = \qty{7.30e-4}{mol/L}$: 73\,\%
extraído.

\textbf{10.} $y = \num{1.0e4}(\num{1.0e-3} - y)(\num{1.0e-2} - y)$ dá $y = \qty{9.89e-4}{mol/L}$: 99\,\%.

\textbf{11.} $y = \num{1.0e4}(\num{1.0e-3} - y)(\num{1.0e-4} - y)$ dá $y = \qty{9.0e-5}{mol/L}$: 9\,\%, limitado pelo \omterm{def:b3:supramolecular:hosts}{éter-coroa},
que está 90\,\% carregado.

\textbf{12.} O íon permanganato conserva sua cor violeta no tolueno. Ali,
ele não está solvatado, é um ânion nu, e reage muito mais depressa que na água, onde uma
camada de hidratação o protege.

\textbf{13.} O íon sódio entra e sai muito mais depressa que a diferença entre as
duas frequências de ressonância: troca rápida.

\textbf{14.} $\delta = \delta_{\mathrm{free}} + (\delta_{\mathrm{bound}} - \delta_{\mathrm{free}})[\mathrm{HG}]/H_0$, com $[\mathrm{HG}]$ dada pela isoterma exata, $H_0 = \qty{2.0e-3}{mol/L}$,
$G_0 = (\text{equiv.})\,H_0$, $\delta_{\mathrm{free}} = \qty{3.560}{ppm}$.

\textbf{15.} $K = (1.55 \pm 0.08) \times 10^3\ \unit{L.mol^{-1}}$, $\delta_{\mathrm{bound}} = \qty{3.719 \pm 0.001}{ppm}$.

\textbf{16.} O resíduo quadrático médio é de \qty{0.0014}{ppm}, da ordem da precisão de
leitura, sem tendência: o modelo 1:1 se ajusta.

\textbf{17.} Fração ligada $f$ em $G_0 = fH_0 + f/(K(1 - f))$: 0.28 equivalente para $f = 0.2$ e 2.1 para $f = 0.8$.

\textbf{18.} $KH_0 = 10^{6.1} \times \num{2.0e-3} \approx 2500$: a curva teria uma quebra nítida em um equivalente
e daria apenas um limite inferior; é necessário um experimento de competição ou uma calorimetria a
concentração mais baixa.

\textbf{19.} No tolueno, uma solução homogênea de alqueno e permanganato.

\textbf{20.} O manganês reduzido deixa a fase orgânica (como dióxido de manganês),
e o \omterm{def:b3:supramolecular:hosts}{éter-coroa}, liberado na interface, capta outro $\ce{K+}$ e outro $\ce{MnO4-}$. Cada molécula de
\omterm{def:b3:supramolecular:hosts}{éter-coroa} transporta o permanganato muitas vezes: uma quantidade catalítica basta.

\textbf{21.} $K_{\mathrm{ex}}$ multiplica $[\ce{K+}]$: uma concentração alta de potássio favorece a extração.

\textbf{22.} Os sais de amônio quaternário, como o cloreto de tetrabutilamônio ou de
benziltrietilamônio, cujo cátion é lipofílico por si só.

\textbf{23.} O benzeno provoca leucemia; o tolueno faz o mesmo trabalho com um perigo muito
menor.

\textbf{24.} Com \qty{1.0}{mmol/L} de \omterm{def:b3:supramolecular:hosts}{éter-coroa}, 73\,\% do permanganato é extraído para o
tolueno.
\end{solution}
